\chapter{Lung Transplantation}
\section*{Overview}
Lung transplantation replaces one or both diseased lungs with donor lungs.
\section*{Common Instruments and Equipment}
The exact setup varies by technique and setting. Common items may include:
\begin{itemize}[nosep,leftmargin=*]
  \item thoracic transplant surgery instrument set
  \item lung preservation and transport equipment
  \item bronchoscope and chest drainage system
\end{itemize}

\section*{Why It Is Done}
It is considered for selected patients with advanced, life-limiting lung disease despite maximal treatment, such as cystic fibrosis, pulmonary fibrosis, or severe pulmonary hypertension.
\section*{How It Is Performed}
The diseased lung or lungs are removed and donor airways and blood vessels are connected under general anesthesia, often with cardiopulmonary support.
\section*{Contraindications}
No single contraindication applies to every patient; the treating team assesses whether the procedure is suitable for the individual and the reason for treatment.

\section*{Precautions}
The care team reviews relevant health conditions, medicines, allergies, and procedure-specific risks, and explains any changes needed for the individual.

\section*{Risks and Follow-up}
Primary graft dysfunction, rejection, infection, airway complications, kidney injury, and medication toxicity are possible. Lifelong immunosuppression and transplant-center monitoring are required.
\section*{References}
\begin{enumerate}\item MedlinePlus. Lung Transplant. U.S. National Library of Medicine; 2025.\item Current Medical Diagnosis and Treatment. McGraw-Hill; 2025.\end{enumerate}
\clearpage
